% Teaching manuscript. Edit this TeX source; Markdown is compatibility navigation only.

\StudyChapter{CS5222-Assignment01}{Assignment 1 · Chapters 1 and 2｜计算、事务与技术证据}
\StudyTarget{CS5222-Assignment01__assignment-1-chapters-1-and-2}


\StudyRepoRef{CS5222/course-notes/README.md}{}{课程目录} · \StudyRef{CS5222-Chapter01}{root}{Chapter1} · \StudyRef{CS5222-Chapter02}{root}{Chapter2} · 原任务（4页）


本讲义提供P6、P31、P17的纸笔推导及Zoom资料核查方法。当前未提供教师解答，以下推导与资料分析供学习参考。截止日期、提交格式和手续回查原任务及\href{https://crazyshout.github.io/micro-course/notices.html}{课程信息栏}。


基础可选：\StudyRef{Foundations-NetworkBasics}{root}{单位与时间轴}、\StudyRef{Foundations-NetworkBasics}{messages}{报文字段}。当前四题都属于核心任务；P6难点是首/末位事件，P31是跨链路重叠，P17是跨会话状态，Zoom是产品情境与来源边界。理解难度均为中等。







\StudyHeading{1}{P6. Transmission and propagation｜最后一位刚出发，第一位可能已经到达}
\StudyTarget{CS5222-Assignment01__p6}
\StudyTarget{CS5222-Assignment01__p6-transmission-and-propagation}


\StudyImage{CS5222/course-notes/assets/assignment01-01.png}


\textbf{条件 / Conditions:} 主机A、B相距m米，链路发送速率R bit/s，传播速度s m/s，包长L bit。忽略处理、排队，从A开始发送首位计时t=0。 / Hosts A and B are m metres apart; link rate is R bit/s, propagation speed s m/s, and packet length L bits. Ignore processing and queueing and start timing with the first bit.


\begingroup
\small\setlength{\tabcolsep}{3pt}
\renewcommand{\arraystretch}{1.22}
\begin{longtable}{@{}>{\raggedright\arraybackslash}p{\dimexpr0.226094\linewidth-1.33333\tabcolsep\relax}>{\raggedright\arraybackslash}p{\dimexpr0.351511\linewidth-1.33333\tabcolsep\relax}>{\raggedright\arraybackslash}p{\dimexpr0.422396\linewidth-1.33333\tabcolsep\relax}@{}}
\toprule
原小问 & 中文推导与答案 & English answer \\
\midrule\endfirsthead
原小问 & 中文推导与答案 & English answer \\
\midrule\endhead
\bottomrule\endfoot
a 传播 & \(d_{prop}=m/s\)，距离除速度 & Propagation delay is m/s. \\
b 发送 & \(d_{trans}=L/R\)，数据量除速率 & Serialization time is L/R. \\
c 整包端到端 & \(L/R+m/s\)，末位先被送入链路再传播 & Complete-packet delay is L/R+m/s. \\
d t=dtrans，末位位置 & 在连续比特模型的边界时刻刚离开A进入链路，还没有走完整条路 & The last bit is just entering the link at A. \\
e dprop>dtrans，首位位置 & 仍在线路内，离A的距离为\(sL/R\lt m\) & The first bit remains on the link, sL/R metres from A. \\
f dprop<dtrans，首位位置 & 已到达B；不能继续说还在链路上 & The first bit has already reached B. \\
g 两时延相等 & \(m/s=L/R\)，故\(m=sL/R\) & Set m=sL/R. \\
\end{longtable}
\endgroup


g给s=2.5\ensuremath{\times}10\textsuperscript{8}m/s、L=120bit、R=56kbps=56000bit/s：
\[
d_{trans}=120/56000\approx0.002142857\text{ s},\quad m\approx535714.29\text{ m}=535.714\text{ km}.\tag{A1.1}
\]


\StudyImage{learning/foundation-notes/assets/network-timeline.png}


图是帮助读事件的教学示意（参数2ms/8ms），不是g题的参数图。g的数值单独复算，避免把示意图标注抄进答案。


\textbf{变式 / Transfer:} L=240bit、R不变，要保持两时延相等，m如何变化？ / Double packet length while keeping R and s fixed; how must m change?



\begin{StudySolution}{答案 / Answer}


m翻倍，约1071.429km。 / m doubles to about 1071.429 km. 改变的是等时延所需距离，不是信号速度。


\end{StudySolution}






\StudyHeading{1}{P31. Message segmentation｜像流水线一样，让不同链路同时工作}
\StudyTarget{CS5222-Assignment01__p31}
\StudyTarget{CS5222-Assignment01__p31-message-segmentation}


\StudyImage{CS5222/course-notes/assets/assignment01-02.png}


\textbf{条件 / Conditions:} 原图源\ensuremath{\to}交换机1\ensuremath{\to}交换机2\ensuremath{\to}目的，共3条2Mbps链路；消息8\ensuremath{\times}10\textsuperscript{6}bit，存储转发，忽略传播、排队和处理。分段时为800个10000bit包，题目未另加包头。 / An 8-million-bit message crosses three 2-Mbps store-and-forward links. Ignore propagation, queueing and processing. Segmentation creates 800 packets of 10000 bits each, with no additional header overhead specified.


a 不分段，第一链路 \(8\times10^6/(2\times10^6)=4\)s；交换机必须等整条消息，3跳总\textbf{12s}。 / Without segmentation: 4 s to the first switch and 12 s to the destination.


b 单包发送10000/(2\ensuremath{\times}10\textsuperscript{6})=\textbf{0.005s=5ms}。第一包到第一交换机t=5ms；它继续去第二交换机时，源发送第二包，第二包在t=\textbf{10ms}完全到第一交换机。 / First packet reaches switch1 at 5 ms; packet2 is completely received there at 10 ms.


c 第一包到目的需3\ensuremath{\times}5ms，后续799包每隔5ms到，总 \((800+3-1)\times0.005=\mathbf{4.01s}\)。 / Completion time is 4.01 s. 链路本身没变快，减少的是“大消息在每一跳全等完”的串行等待。


\StudyImage{CS5222/course-notes/assets/segmentation.png}


图中展开前几个包，末尾标出全部800包的完成时刻。


d 其他好处：较短包给其他流量穿插机会；出现丢失时可只重传受影响部分（若上层协议支持）；较小存储/处理单位便于逐包转发。 / Segmentation enables interleaving, selective retransmission when supported, and manageable forwarding units.


e 代价：每包头部和处理开销，更多调度；可能乱序，需要识别与重组；重传/超时机制增加复杂度。 / Costs include headers, per-packet processing, possible reordering, reassembly and recovery complexity. 不把“越小永远越快”当结论；本题忽略头部是有意简化。


\textbf{变式 / Transfer:} 同消息改成400包，每包20000bit，其他不变，总多久？ / Use 400 packets of 20000 bits with the same path.



\begin{StudySolution}{答案 / Answer}


每包10ms，总(400+3\ensuremath{-}1)\ensuremath{\times}0.01=4.02s。 / 4.02 s. 仍近似受源端总上传4s约束，只是流水线填充稍慢。


\end{StudySolution}






\StudyHeading{1}{P17. POP3｜完成整次事务，不只补两行}
\StudyTarget{CS5222-Assignment01__pop3}
\StudyTarget{CS5222-Assignment01__p17-pop3}


\StudyImage{CS5222/course-notes/assets/assignment01-03.png}


\textbf{完整条件 / Conditions:} 已认证进入transaction，LIST给两封邮件，编号1/2、大小498/912byte；原片段已经RETR1。a为download-and-delete，b为download-and-keep，均继续到服务完成。c在b后退出，5分钟后再次登录，无新邮件，给出第二会话。 / Two messages of 498 and 912 octets are present. Finish the full download-and-delete and download-and-keep transactions, then show a second keep-mode session five minutes later with no new mail.


原题略去若干\StudyInline{+OK}确认；下面补齐规范响应。\StudyInline{[message .\allowbreak{}.\allowbreak{}.\allowbreak{}]}只是未知邮件内容的说明占位，\textbf{不是实际发送的协议文本}。每行以CRLF结束，单独\StudyInline{.\allowbreak{}}结束多行结果。命令字大小写不影响本题含义。



\StudyHeading{2}{a. Download and delete / 下载并删除}
\StudyTarget{CS5222-Assignment01__a-download-and-delete}



\begin{StudyCode}
C: LIST
S: +OK 2 messages (1410 octets)
S: 1 498
S: 2 912
S: .
C: RETR 1
S: +OK 498 octets
S: [message 1 content]
S: .
C: DELE 1
S: +OK message 1 marked for deletion
C: RETR 2
S: +OK 912 octets
S: [message 2 content]
S: .
C: DELE 2
S: +OK message 2 marked for deletion
C: QUIT
S: +OK signing off
\end{StudyCode}


DELE先标记，成功QUIT使会话进入UPDATE并提交删除；本例假设操作成功且没有异常退出。只写DELE1然后QUIT会漏掉第二封的下载与删除。 / Retrieve both messages, mark both for deletion, then successfully quit to commit the update.



\StudyHeading{2}{b. Download and keep / 下载并保留}
\StudyTarget{CS5222-Assignment01__b-download-and-keep}



\begin{StudyCode}
C: LIST
S: +OK 2 messages (1410 octets)
S: 1 498
S: 2 912
S: .
C: RETR 1
S: +OK 498 octets
S: [message 1 content]
S: .
C: RETR 2
S: +OK 912 octets
S: [message 2 content]
S: .
C: QUIT
S: +OK signing off
\end{StudyCode}


没有DELE，两封保留。RETR不会自动删除；“keep”不是另加一个名叫KEEP的POP3命令。 / No messages are marked for deletion, so both remain on the server.



\StudyHeading{2}{c. Second keep-mode session / 五分钟后的第二次会话}
\StudyTarget{CS5222-Assignment01__c-second-keep-mode-session}


假设没有额外客户端UIDL去重、服务器未改变这两封邮件的编号，重新认证后仍LIST两封；可以再次RETR1/2。原问题仅说“无新邮件”，没有说服务器将已读邮件隐藏。



\begin{StudyCode}
S: +OK POP3 server ready
C: USER <username>
S: +OK
C: PASS <password>
S: +OK authenticated
C: LIST
S: +OK 2 messages (1410 octets)
S: 1 498
S: 2 912
S: .
C: RETR 1
S: +OK 498 octets
S: [message 1 content]
S: .
C: RETR 2
S: +OK 912 octets
S: [message 2 content]
S: .
C: QUIT
S: +OK signing off
\end{StudyCode}


尖括号内为用户名和密码的占位。实际客户端可保存UIDL跳过已下载邮件；若采用这种扩展约定，可能LIST/UIDL核对后不RETR旧邮件。应明确这是客户端保存状态的行为，而不是“POP3服务器因为上次读过就删掉”。\href{https://www.rfc-editor.org/rfc/rfc1939}{RFC1939}（查阅于2026-09-23）


\textbf{English answer:} With no deletion and no client-side duplicate suppression, both messages are still available and can be downloaded again. State any UIDL-based client behavior explicitly.


\textbf{变式 / Transfer:} b中仅DELE1并成功QUIT，之后无新邮件，剩多少总字节？ / Delete only message1 in part b and successfully quit. How many message octets remain?



\begin{StudySolution}{答案 / Answer}


只剩原邮件2，912byte；新会话编号不应依赖旧编号恒定。 / 912 octets remain; message numbers need not remain identical across sessions.


\end{StudySolution}






\StudyHeading{1}{Zoom. Apply protocol ideas with sources｜先说研究的是哪一种Zoom连接}
\StudyTarget{CS5222-Assignment01__zoom}
\begingroup
\titlespacing*{\subsection}{0pt}{.7\baselineskip}{.25\baselineskip}
\titlespacing*{\section}{0pt}{1.1\baselineskip}{.4\baselineskip}

\StudyTarget{CS5222-Assignment01__zoom-apply-protocol-ideas-with-sourceszoom}


\StudyImage{CS5222/course-notes/assets/assignment01-04.png}


\textbf{原问题 / Questions:} a列举三种应用层协议并解释作用；b主要client-server还是P2P；c主要传输协议及原因；d学生是否需要老师电脑IP及原因。 / Identify three application-layer protocols and their roles, explain the main architecture and transport choice, and determine whether a student needs the lecturer's host IP.


\textbf{核查情境：2026-09-23，普通多人Zoom Meetings、原生Zoom Workplace客户端、经Zoom云入会的默认讨论范围。} 没有实际抓取本校课堂流量，也未确认本校启用哪些企业功能。因此下述为官方资料支撑的典型架构分析，不是对某一次课堂连接的实测报告。



\StudyHeading{2}{a. 应用协议：产品生态的三项，与原生媒体路径分开}
\StudyTarget{CS5222-Assignment01__a}


题目未限定客户端类型，可用下表列出Zoom产品生态明确支持的三项，并把各自情境写进答案：


\begingroup
\small\setlength{\tabcolsep}{3pt}
\renewcommand{\arraystretch}{1.22}
\begin{longtable}{@{}>{\raggedright\arraybackslash}p{\dimexpr0.109893\linewidth-1.50000\tabcolsep\relax}>{\raggedright\arraybackslash}p{\dimexpr0.213670\linewidth-1.50000\tabcolsep\relax}>{\raggedright\arraybackslash}p{\dimexpr0.388840\linewidth-1.50000\tabcolsep\relax}>{\raggedright\arraybackslash}p{\dimexpr0.287598\linewidth-1.50000\tabcolsep\relax}@{}}
\toprule
协议 & 功能（中文） & English explanation & 证据/适用情境 \\
\midrule\endfirsthead
协议 & 功能（中文） & English explanation & 证据/适用情境 \\
\midrule\endhead
\bottomrule\endfoot
HTTPS & 网页/客户端与Zoom Web服务的请求、资源获取等 & Secure HTTP access to web services and content & \href{https://support.zoom.com/hc/en/article?id=zm\_kb\&sysparm\_article=KB0060548}{官方网络配置}明确HTTP/HTTPS访问与CDN；不据端口猜全部音视频都经HTTP \\
SIP & 让标准会议室端点进行会话建立/互通 & Session signaling for supported standards-based room endpoints & \href{https://support.zoom.com/hc/en/article?id=zm\_kb\&sysparm\_article=KB0060661}{Conference Room Connector}；不是所有原生客户端必经路径 \\
H.323（协议族） & 另一类标准音视频会议端点经连接器接入 & A conferencing protocol suite for supported room-system interoperability & 同一官方连接器资料；需注明它是协议族而非单一报文协议 \\
\end{longtable}
\endgroup


表中HTTPS用于Web服务，SIP/H.323用于会议室端点互通，作答时分别注明情境。若教师将题意限定为\textbf{原生客户端自身的媒体/控制三种协议}，应按该范围补充特定版本的技术资料，而不能把SIP/H.323互通项代替原生客户端路径的说明。公开端口表只能证明传输端口用途，不能单独证明RTP/SRTP等内部封装；作答时应注明所选客户端与部署情境。



\StudyHeading{2}{b. 架构 / Architecture}
\StudyTarget{CS5222-Assignment01__b-architecture}


普通多人云会议主要由云服务协调与转发，按课程分类主要是client-server：各客户端向服务基础设施连接，云处理参与者发现、会话与媒体分发。\textbf{这是基于官方云连接/回退说明的架构归纳}，具体部署也可能包含下面介绍的P2P或Mesh连接。\href{https://library.zoom.com/advanced-enterprise-services/zoom-mesh/zoom-mesh-explainer}{Zoom Mesh官方说明}


\textbf{English:} Ordinary multi-participant cloud meetings are primarily client-server in the course's sense; clients rely on Zoom infrastructure for coordination and media distribution.


明确例外：官方允许同一局域网的两人会议在符合设置条件时用peer-to-peer；加入第三人等情况会使用云。\href{https://support.zoom.com/hc/en/article?id=zm\_kb\&sysparm\_article=KB0067877}{两人P2P说明} 企业启用Zoom Mesh时，局域网内可由parent客户端给child分发视频；这是有云协调的混合机制，不能拿来证明普通多人课堂始终直接连接老师。



\StudyHeading{2}{c. 传输 / Transport}
\StudyTarget{CS5222-Assignment01__c-transport}


实时媒体通常优先UDP，官方DSCP文档说明默认媒体在UDP8801等端口，特定管理设置可拆分音频/视频/共享端口。\href{https://library.zoom.com/admin-corner/network-management/quality-of-service-and-network-best-practices-explainer/using-dscp-marking-with-zoom}{官方媒体端口说明} 控制/Web访问及媒体回退可以使用TCP；“Zoom只用UDP”错误。课程层面的原因是实时音视频更在意及时性，可由应用处理丢失与适应码率，而TCP有序交付可能让后续数据等重传；这解释选择动机，不意味UDP本身自动更可靠或保证低延迟。


\textbf{English:} Real-time media commonly prefers UDP, while control/web traffic and fallback paths may use TCP. Low delay and application-level adaptation matter more than waiting indefinitely for every media byte.



\StudyHeading{2}{d. 老师电脑的IP / Lecturer's host IP}
\StudyTarget{CS5222-Assignment01__d-ip-lecturers-host-ip}


在普通云课堂情境，学生连接的是Zoom提供的服务端点，不需要自己先解析或输入老师个人电脑的IP。老师和学生各自向基础设施建立通信，服务负责转发；知道meeting ID也不是把它翻译成老师IP。


\textbf{English:} In an ordinary cloud meeting, the student does not need to obtain the lecturer's host IP; Zoom infrastructure routes the session. Peer/mesh configurations may use other local endpoints and require separate analysis.



\StudyHeading{2}{怎样整理成自己的作答}
\StudyTarget{CS5222-Assignment01___1}


每个判断写成“结论 \ensuremath{\to} 明确情境 \ensuremath{\to} 机制 \ensuremath{\to} 官方来源及查阅日期”。区分产品支持、当前默认与本机实测三种证据。原题是应用课程知识的调查题，不应把一段不知年份的博客当软件永久架构。



\StudyHeading{1}{覆盖、复算与验收}
\StudyTarget{CS5222-Assignment01___2}


原PDF p1的P6(a–g)、p2的P31(a–e)及原拓扑、p3的P17(a–c)及“写到事务结束”要求、p4的Zoom(a–d)全部对应。数值和分包图由 \StudyRepoRef{CS5222/course-notes/tools/network_examples.py}{}{network\_examples.py} 复算；原题未给邮件正文，事务中使用占位说明。


完成前检查：每个速率/长度有单位，每段延迟不多算漏算；POP3包含两封邮件与成功结束；Zoom区分客户端/部署并注明来源时间。再回到\href{https://crazyshout.github.io/micro-course/cards.html\#CS5222-N260}{已有N260–N263}复习对应知识。
\endgroup
